\section{Attività di revisione scientifica}

\outerlist{

\entrybig
	{\textbf{Guest Editor}}{}
	{\textcolor{darkgray}{Special Issue su rivista internazionale Q1 secondo Ranking SCImago 2023} [\href{https://www.mdpi.com/journal/pathogens/special_issues/LORRUF04K5}{link}]}{2022}
\innerlist{
    \entry{\textbf{Attività:} Guest Editor dello \href{https://www.mdpi.com/journal/pathogens/special_issues/LORRUF04K5}{\textit{Special Issue} “Artificial Intelligence (AI) as One Health Tool in Pathology to Monitor and Control Infectious Diseases”} della rivista \textit{Pathogens}.}
}

\entrybig
	{\textbf{Revisore di tesi di dottorato}}{2025 -- in corso}
	{\textcolor{darkgray}{Programmi di Dottorato nazionali e internazionali}}{}
\innerlist{
    \entry{Tomaso Trinci (2025), Università degli Studi di Firenze.}
    \entry{Federico Girella (2026), Università degli Studi di Verona.}
    \entry{Amelia Sorrenti (2026), National PhD Program in Artificial Intelligence.}
}

\entrybig
	{\textbf{Revisore di progetti di ricerca post-dottorali}}{}
	{\textcolor{darkgray}{Research Foundation Flanders (FWO) -- European Science Foundation}}{2025}
\innerlist{
    \entry{\textbf{Attività:} valutazione come esperto esterno di una proposta di borsa post-dottorale nell’ambito del bando \textit{Junior and Senior Postdoctoral Fellowship 2024} della Research Foundation Flanders (FWO), con processo di revisione indipendente gestito dalla European Science Foundation (ESF).}
    \entry{\textbf{Progetto valutato:} \textit{“Making Continual Learning Ready for the Dynamic Real-world: Long-term Unsupervised Continual Representation Learning”} (Ref.~24-FWO-PDOC-1509).}
}

\entrybig
	{\textbf{Revisore scientifico}}{2019 -- in corso}
	{\textcolor{darkgray}{Conferenze internazionali top-tier in AI, ML e Computer Vision}}{}
\innerlist{
    \entry{\textbf{Attività:} attività di revisione per le principali conferenze internazionali nei settori della \textit{Computer Vision}, del \textit{Machine Learning} e dell’\textit{Artificial Intelligence}, tra cui CVPR, ICCV, ECCV, NeurIPS, IJCAI e ICPR.}
}
}